% =====================================================================
% COVER LETTER — Radiology: Artificial Intelligence
%
% Uploaded as its own file. NOT anonymized: the cover letter goes to the
% editorial office, not to peer reviewers, and it is the only place the real
% repository URL and the author identities may appear.
%
% RSNA specifies six items for the cover letter, four of them flagged REQUIRED
% (marked below). Two further requirements come from RSNA Editorial Policies:
% the AI-use disclosure and the data-control statement.
%
% Deliberately NOT here, per the instructions:
%   - suggested / opposed reviewers  -> entered in ScholarOne Step 5
%   - a "why this journal" pitch     -> only a fast-track request takes one,
%                                       and that must also be emailed to
%                                       radiology-ai@rsna.org
%
% Everything in [SQUARE BRACKETS] is an author action.
% =====================================================================
\documentclass[11pt]{article}
\usepackage[margin=1in]{geometry}
\usepackage[hidelinks]{hyperref}
\usepackage{parskip}
\usepackage{times}
\pagestyle{empty}

\begin{document}

[Author name, degrees] \\
{}[Department], [Institution] \\
{}[Address] \\
{}[e-mail] \quad [telephone]

\vspace{1.2em}
[Date]

\vspace{1.2em}
Charles E. Kahn, Jr., MD, MS \\
Editor, \textit{Radiology: Artificial Intelligence} \\
Radiological Society of North America

\vspace{1.2em}
Dear Dr Kahn,

\textbf{Title.} ``Supervised Tumor Attention Transfers Across CT Cohorts, but
\textit{EGFR} Prediction Does Not Exceed Clinical Variables.''

\textbf{Authors.} [Complete author list, in final order, with degrees.]

\textbf{Article type.} Original Research.

\medskip
We submit this manuscript for your consideration. The CT radiogenomics
literature reports \textit{EGFR} prediction at AUC 0.80--0.89, almost always from
single-center cohorts without external validation. We took an image encoder whose
attention is supervised against tumor segmentations and separated two claims this
literature routinely bundles: that such a model attends to the lesion, and that it
genotypes it. The attention transferred to 420 patients in another country, on
another scanner fleet, in another treatment setting. Genotype prediction did not: pooled AUC was 0.618 against 0.794 for smoking status
alone, adding the model to a five-variable clinical model changed AUC by
$-0.029$, no fusion topology helped, and calibration was worse than predicting the
base rate.

We think the useful result is the dissociation rather than the null. A saliency
map that tracks the tumor on an unseen institution's data, attached to a
classifier performing at 0.618, is evidence that plausible localization is not
evidence of predictive validity --- and that can only be observed when the
prediction fails, so it cannot appear in studies that report both as successes.

\textbf{Two objections we would rather raise than have noticed.} First, the
internal cohort is 153 patients with a known \textit{EGFR} call and the
pre-specified primary analysis set is 97, far below what the journal typically
carries. Our argument is not that this is adequate but that the question is not
resolvable by enlarging it: Riley's criterion requires 414 patients for the
five-variable clinical model alone, the entire open world of
\textit{EGFR}-labeled CT gives a 95\% confidence-interval half-width of 0.154, and
the largest external study yet published (1646 patients, 11\,473 lesions,
\textit{Eur Radiol} 2026) reproduces our interval at ten times the sample size
with our internal cohort as its external one. Second, the primary endpoint clears
its pre-specified threshold by a single fold (11 of 15). Table 3 therefore reports
every evaluation individually; the distribution is bimodal rather than marginal,
with the weakest success at 0.757 pointing accuracy and the strongest failure at
exactly 0.000.

\textbf{Two things a reviewer would ask for, done in advance.} The crops are
lesion-centered, so a center prior is a serious control rather than a courtesy. We
therefore translated each external crop and its reference contour so the tumor was
no longer centered: across the aligning folds pointing accuracy fell from 0.831 to
0.711 at a 35\% displacement, while the center prior fell from 0.476 to 0.002.
Separately, because 15 fold-checkpoints on one external cohort are not 15
independent studies, we reran the endpoint as a patient-level bootstrap. It
favours the model when the training runs are held fixed and does \textit{not}
survive treating the run as a random effect. We report that against our own
interest, and it is the paper's main limitation.

\medskip
\textbf{Subject overlap with prior work} \textit{(required)}. Both imaging
collections are public and have been analyzed by other groups, including in each
collection's source publication and in the multi-tumor external-validation study
cited above, in which NSCLC-Radiogenomics serves as the external cohort. No
patient in this analysis overlaps with any published, in-press, or under-review
work by these authors. What is new here is the pre-specified separation of the
localization and genotyping claims, the three-control external localization
evaluation, and the measured precision bound. This overlap is also reported in
Materials and Methods. [If any overlapping article exists, upload its PDF at
submission and describe here what is new and different.]

\textbf{Conflicts of interest and industry support} \textit{(required)}.
[REQUIRED --- state all conflicts and any industry support, or state that there
are none. ICMJE disclosure forms are submitted for every author.]

\textbf{Sole submission} \textit{(required)}. This manuscript is original, has not
been published in the same or similar form by the authors, and has not been and
will not be submitted to another journal by the authors or by colleagues at their
institutions while it is under consideration at \textit{Radiology: Artificial
Intelligence}. [State here if any similar material has been provided to news or
advertising media, and supply a copy.]

\textbf{Use of artificial intelligence in the study and in manuscript
preparation} \textit{(required by RSNA Editorial Policies)}. Generative AI was
used substantially in this work, and we would rather state the extent than have it
inferred.

\textit{Tool.} Claude Opus 5, accessed through the Claude Code command-line
interface (Anthropic, San Francisco, Calif). \textit{Dates of access.} July 20 to
August 13, 2026. \textit{Use.} Writing and refactoring the analysis code;
executing and summarising the reported analyses; literature search; drafting and
revising the manuscript, the figures' generating code, and the reporting
checklist. Sixty-five of the 68 commits in the repository record that assistance
in their commit metadata, so the claim is auditable rather than asserted.

\textit{What the authors did.} The authors defined the research question,
approved the pre-specified analysis plan before external evaluation, verified
every reported number against the archived artifact from which it is regenerated,
and take full responsibility for the integrity and accuracy of the work. No AI
tool meets ICMJE authorship criteria and none is listed as an author.

\textit{Errors the assistance introduced.} We think a disclosure that reports no
errors is describing its review process rather than its work, so: the assistance
applied a published sample-size formula incorrectly and the wrong minimum-$n$
reached a draft; it mis-attributed a cited paper, describing a
blood-brain-barrier reconstruction study as a mutation-prediction study; and a
stale command-line default caused the fusion comparison to be computed on a
variant model rather than the reported one, producing an internally inconsistent
table. Each was caught before submission, each is recorded with its correction in
the public version history, and each has a regression test written in response.
The last of these was found by an adversarial review of the manuscript, which is
also disclosed here because it materially improved the paper.

\textbf{Control of the data and the analyses} \textit{(required by RSNA Editorial
Policies)}. [REQUIRED --- name the author(s) who had full access to and control of
the data and who performed the statistical analyses, e.g.\ ``[Initials] had full
access to all study data and takes responsibility for the integrity of the data
and the accuracy of the analysis.'']

\textbf{Dual first authorship.} [Delete if not applicable. Otherwise state which
two authors contributed equally and why equal contribution is warranted; this must
also appear on the full title page.]

\textbf{Prior presentation and preprint posting.} [REQUIRED disclosure at
submission --- state any presentation at a scientific meeting and any preprint
posting with server, DOI and date, or state that there has been none.]

\medskip
\textbf{Reproducibility, for the editorial office.} Every number in the
manuscript is regenerated from archived out-of-fold predictions by a single
script, requiring no imaging data and no GPU. The repository holds the
pre-specified analysis plan with its deviation log --- including the deviations
that went against us --- the completed CLAIM checklist, and all archived runs,
negative results included. Because review is double-anonymized, the manuscript
cites an anonymized mirror; the permanent URL is \texttt{[REAL REPOSITORY URL]}
and the archival DOI is \texttt{[ZENODO DOI]}, given here for the editorial office
only.

\textbf{Image licensing.} NSCLC-Radiogenomics is licensed CC BY 3.0, and its
images appear in Figure 6 with attribution. NSCLC-Radiomics is licensed CC BY-NC
3.0, so no pixel data from it appears in any figure --- only derived statistics.

\textbf{Ethics.} [REQUIRED --- insert the institutional review board's written
determination. The analysis used only publicly available, fully de-identified data
distributed by The Cancer Imaging Archive under open licenses, and no protected
health information was accessed.]

\vspace{1.2em}
Yours sincerely,

\vspace{2.5em}
{}[Corresponding author name, degrees] \\
on behalf of all authors

\end{document}
